% !TeX root = ../../artigo-exemplo-undf.tex
\undfsetcontent{resumo}{%
A reutilização de uma classe LaTeX em diferentes modalidades de trabalho
acadêmico exige verificar numeração, bibliografia e continuidade do texto.
Este estudo didático tem como objetivo examinar a consistência estrutural dos
perfis de monografia e artigo do UnDF-TeX. O procedimento consiste na análise
de casos mínimos com dados sintéticos, compilação por LuaLaTeX e Biber,
extração de texto e inspeção dos documentos produzidos. O corpus compreende
quatro modelos principais e casos específicos de complementos, citação de
citação, objetos extensos e inclusão de PDF. Na revisão local de 14 de outubro
de 2026, os 48 destinos de referências cruzadas dos complementos corresponderam
às páginas registradas nos arquivos auxiliares. Os quatro cenários de citação
de citação produziram as listas bibliográficas previstas, e os testes de
objetos extensos e inclusão de PDF passaram. Os resultados sustentam a
consistência dos comportamentos examinados na versão observada. Não permitem
concluir conformidade normativa integral, aprovação institucional ou ausência
de defeitos em documentos arbitrários. Conclui-se que casos mínimos e revisão
visual oferecem evidências complementares para a manutenção do modelo.}
\undfsetcontent{palavraschave}{documentos acadêmicos; LaTeX; testes de regressão; referências cruzadas}
